\documentclass[10pt,a4paper]{amsart}
\usepackage[margin=19mm]{geometry}
\usepackage{amsmath,amssymb,mathtools}
\usepackage[scr=rsfs,cal=euler]{mathalfa}
\usepackage{xfrac}
\usepackage[unicode,hidelinks,bookmarks=false]{hyperref}
\newcommand{\ie}{i.e.\ }
\newcommand{\cf}{cf.\ }
\newcommand{\resp}{respectively\ }
\newcommand{\Subsection}{\subsection}
\providecommand{\nobreakdash}{\nobreak\hbox{-}}
\setlength{\emergencystretch}{2em}
\multlinegap0pt
\usepackage{amssymb}
\usepackage{xcolor}
\usepackage{float}
\newtheorem{proposition}{Proposition}
\newtheorem{lemma}{Lemma}
\newtheorem{remark}{Remark}
\def\FF{\ensuremath{\mathcal F}}
\def\R{\ensuremath{\mathbb{R}}}
\def\TT{\ensuremath{\mathcal T}}
\newcommand{\into}{\hookrightarrow}
\newcommand{\onto}{\twoheadrightarrow}
\title{Rank Two Sheaves With Maximal Third Chern Class on the Fano Threefold of Index 2 and Degree 5}
\author{Danil A. Vassiliev}
\date{Source: arXiv:2608.26928v1 (CC BY 4.0)}
\begin{document}
\maketitle

\noindent\emph{Source and licence.} Rank Two Sheaves With Maximal Third Chern Class on the Fano Threefold of Index 2 and Degree 5. Version arXiv:2608.26928v1 (CC BY 4.0), distributed by arXiv under the Creative Commons Attribution 4.0 International licence, http://creativecommons.org/licenses/by/4.0/, as recorded in the arXiv licence metadata for this exact version and shown on the abstract page https://arxiv.org/abs/2608.26928v1. Full licence notice: \texttt{licenses/arxiv-2608.26928-CC-BY-4.0.txt}.\par\smallskip

\noindent\textit{Editorial scope.} Counted unit: the article's lemma -3/5 (source0.tex lines 676--680). The article's complete original proof of it is delivered with it (source0.tex lines 682--699, one \texttt{proof} environment of 18 lines). No mathematics is written, repaired or re-derived: the statement and the proof are the article's own lines, in the article's own notation, and no retained line is substituted or re-flowed.  Retained uncounted as the source-local context the proof needs: lines 340--349; lines 410--414; lines 429--436; lines 440--442; lines 444--448; lines 649--650; lines 655--656.  Nothing was omitted from any retained span, so the package carries no omission record.  No retained line contains a citation, so the wrapper carries no bibliography and no bibliography entry.  No retained line contains a figure environment, an image-inclusion command or drawing code, so no image is delivered and the package declares no asset dependency.  Editorial formatting only, in addition: an AMS \texttt{amsart} preamble that restates the article's own declarations and macros used by the retained lines, namely the environments \texttt{proposition}, \texttt{lemma}, \texttt{remark} on separate counters, together with the standard packages those lines need.  The retained environments print in this document as Proposition 1, Lemma 1, Remark 1 in order of appearance, whereas the article prints the same items with the numbering of its own full document.  The complete native source of the article as distributed in the arXiv e-print for this exact version is bundled unchanged as \texttt{sources/208672/source0.tex}.
	\begin{proposition}[Properties of actual walls]\label{actual}Let $0\to F\to E\to G\to 0$ be an exact sequence in $\mathrm{Coh}^\beta(X)$ defining a wall $W$.
		\begin{enumerate}
			\item If $W$ is semicircular and $\mathrm{ch}_0(F)>\mathrm{ch}_0(E)\ge 0$, then
			$$\rho(E,F)^2\le\frac{\overline{\Delta}_H(E)}{4H^3\cdot\mathrm{ch}_0(F)(H^3\cdot(\mathrm{ch}_0(F)-\mathrm{ch}_0(E))}.$$
			\item If $W$ is semicircular, then $\overline{\Delta}_H(F)+\overline{\Delta}_H(G)<\overline{\Delta}_H(E).$
			\item For any point $(\beta,\alpha)\in W$, we have $0\le\mathrm{ch}_1^\beta(F)\le\mathrm{ch}_1^\beta(E)$. In case of equality on either end, the wall is vertical. In particular, if $H^2\cdot\mathrm{ch}_1^{\beta_0}(E)>0$ has the minimal positive value among objects of $\mathrm{Coh}^{\beta_0}(X)$, then there is no actual wall that intersects the ray $\beta=\beta_0$.
			\item If $\mathrm{ch}_0(E)>0$, then either $\mathrm{ch}_0(F)>0$ and $\mu(F)\le\mu(E)$ or $\mathrm{ch}_0(G)>0$ and $\mu(G)\le\mu(E)$. 
			\item Suppose that $W$ is semicircular. If $\mathrm{ch}_0(E)>0, \mathrm{ch}_0(F)>0$ and $\mu(F)\le\mu(E)$, then $\beta_-(E)<\beta_-(F)\le\mu(F)<\mu(E)$. Respectively, if $\mathrm{ch}_0(E)>0, \mathrm{ch}_0(G)>0$ and $\mu(G)\le\mu(E)$, then $\beta_-(E)<\beta_-(G)\le\mu(G)<\mu(E)$. 
		\end{enumerate}
	\end{proposition}

	\begin{lemma}\label{nu-lambda}
		Let $E\in\mathrm{Coh}^{\beta_0}(X)$ be a $\nu_{\alpha_0,\beta_0}$-semistable object such that $\nu_{\alpha_0,\beta_0}(E)=0$ and fix some $s>0$. Then $\lambda_{{\alpha_0},{\beta_0},s}(E[1])=+\infty$ and $E[1]$ is $\lambda_{{\alpha_0},{\beta_0},s}
		$-semistable. If moreover $E$ is $\nu_{\alpha_0,\beta_0}$-stable, then there is a neighbourhood $U$ of $(\beta_0,\alpha_0)$ such that such that for all $(\beta,\alpha)\in U$ with $\nu_{\alpha,\beta}(E)
		>0$ the object $E$ is $\lambda_{\alpha,\beta,s}$-semistable.
	\end{lemma}

		\begin{equation}\label{T'' F''}
				\begin{split}
						& \TT''_{\gamma}=\{E\in\mathcal A^{\alpha,\beta}(X)\ |\ \text{any quotient}\ E\onto G
						\,\text{satisfies}\,\lambda_{\alpha,\beta, s}(G)>\gamma \}, \\
						& \FF''_{\gamma} = \{E \in \mathcal A^{\alpha, \beta}(X)\ |\ \text{any
								subobject}\ 0\ne F\into E\,\text{satisfies}\, \lambda_{\alpha,\beta, s}(F)\leq\gamma\}.
					\end{split}
			\end{equation}

	\begin{proposition}\label{heart}Let $(\beta,\alpha)\in D$, and let $ \TT''_{\gamma}$, $\FF''_{\gamma}$ be the torsion pair on the category $\mathcal A^{\alpha,\beta}(X)$ defined in (\ref{T'' F''}). There are $s>\frac 16$ and $\gamma\in\R$ such that
		$\langle \mathcal T''_{\gamma},\mathcal F''_{\gamma}[1]\rangle =\mathfrak C.$
	\end{proposition}

	\begin{remark}\label{complexes}
		Any object $E\in\mathfrak C$ is isomorphic to a complex of the form
		$$\mathcal O_{X}(-1)^{\oplus a}\to\mathcal Q(-1)^{\oplus b}\to\mathcal U^{\oplus c}\to\mathcal O_{X}^{\oplus d},$$
		concentrated in degrees $-3,-2,-1$ and $0$. The numbers $a,b,c,d$ are uniquely determined by the Chern character of $E$.
	\end{remark}

	\begin{lemma}\label{3 -2 3/5}There are no tilt-semistable objects $E\in\mathrm{Coh}^\beta(X)$ with $\mathrm{ch}_{\le 2}(E)=3-2H+\frac 35H^2$.
	\end{lemma}

	\begin{lemma}\label{4 -2 2/5} There are no tilt-semistable objects $F\in\mathrm{Coh}^\beta(X)$ with $\mathrm{ch}_{\le 2}(F)=4-2H+\frac 25H^2$.
	\end{lemma}

	\begin{lemma}\label{2 0 -3/5} If $E\in\mathrm{Coh}^\beta(X)$ is tilt-semistable and $\mathrm{ch}(E)=2-\frac 35H^2+eH^3$, then $e\le 0$, and in case of equality $E$ is the cohomology sheaf of a monad
		\begin{equation}\label{3-term}
			\mathcal Q(-1)^{\oplus 3}\to\mathcal U^{\oplus 6}\to\mathcal O_X.
		\end{equation}
	\end{lemma}

	\begin{proof} We have $\beta_-(E)=-\sqrt{\frac 35}\in(-1,-\frac 12)$. The hyperbola $\nu_{\alpha,\beta}(E)=0$ intersects the line $\alpha=\beta+1$ at the point with $\alpha=\frac 15$, hence either $E$ is $\nu_{\alpha,\beta}$-semistable at some point $(\beta,\alpha)\in D$ with $\nu_{\alpha,\beta}(E)=0$ (so $E[1]\in\mathfrak C$ by Lemma \ref{nu-lambda} and Proposition \ref{heart}), or $E$ is tilt-semistable above some semicircular wall of radius $\rho\ge\frac 15$. Assume that there is such a wall, induced by a subobject or quotient $F$ with $\mathrm{ch}_{\le 2}(F)=s+xH+yH^2,s>0,x<0$. If $s>2$, then Proposition \ref{actual}~(1) implies that $\frac{60}{100s(s-2)}\ge\frac{1}{25}$, that is, $s\le 5$. Hence we can assume that $1\le s\le 5$.
		
		We have $0<\mathrm{ch}_1^{\beta_-(E)}(F)<\mathrm{ch}_1^{\beta_-(E)}(E)$, i.e.
		\begin{equation}\label{x4}0<x+\frac{\sqrt{15}}{5}s<\frac{2\sqrt{15}}{5}.
		\end{equation}
		
		If $s=1$, then (\ref{x4}) implies that $x=0$, contradicting our assumptions.
		
		If $s=2$, then (\ref{x4}) gives $x=-1$. We should have $\overline{\Delta}_H(F)=25-100y\in[0,60)$ and $y\in\frac 12+\frac 15\mathbb Z$, hence $y\in\{-\frac{3}{10},-\frac{1}{10},\frac{1}{10}\}$. Since $\overline{\Delta}_H(F)+\overline{\Delta}_H(E/F)=\overline{\Delta}_H(F)+25<\overline{\Delta}_H(E)=60$, only the case $y=\frac{1}{10}$ remains possible. However, in this case the equation defining $W(E,F)$ defines an empty set.
		
		If $s=3$, then (\ref{x4}) implies that $x\in\{-2,-1\}$. Assume that $x=-2$, then $\overline{\Delta}_H(F)=100-150y\in[0,60)$ and $y\in\frac 15\mathbb Z$, hence $y\in\{\frac 25,\frac 35\}$. The condition $\overline{\Delta}_H(F)+\overline{\Delta}_H(E/F)<\overline{\Delta}_H(E)$ rules out the first case, and Lemma \ref{3 -2 3/5} prohibits the second. Assume now that $x=-1$, then $\overline{\Delta}_H(F)=25-150y\in[0,60)$ and $y\in\frac 12+\frac 15\mathbb Z$, hence $y\in\{-\frac{1}{10},\frac{1}{10}\}$. The condition $\overline{\Delta}_H(E/F)\ge 0$ implies that only the case $y=-\frac{1}{10}$ remains possible. However, in this case $W(E,F)$ is on the wrong side of the vertical wall.
		
		If $s=4$, then (\ref{x4}) gives $x\in\{-3,-2\}$. Let $x=-3$, then $\overline{\Delta}_H(F)=225-200y\in[0,60)$ and $y\in\frac 12+\frac 15\mathbb Z$, so $y\in\{\frac{9}{10},\frac{11}{10}\}$. We have $\overline{\Delta}_H(E/F)=165-100y$, and the property $\overline{\Delta}_H(F)+\overline{\Delta}_H(E/F)<\overline{\Delta}_H(E)$ in both cases is not satisfied. Let now $x=-2$, then $\overline{\Delta}_H(F)=100-200y\in[0,60)$ and $y\in\frac 15\mathbb Z$, hence $y=\frac 25$. However, by Lemma \ref{4 -2 2/5} there are no tilt-semistable objects $F$ with such a Chern character.
		
		If $s=5$, then (\ref{x4}) implies that $x=-3$. In this case $\overline{\Delta}_H(F)=225-250y\in[0,60)$ and $y\in\frac 12+\frac 15\mathbb Z$, hence $y\in\{\frac{7}{10},\frac{9}{10}\}$. The condition $\overline{\Delta}_H(F)+\overline{\Delta}_H(E/F)<\overline{\Delta}_H(E)$ is satisfied only for $y=\frac{9}{10}$. However, in this case $W(E,F)$ is empty.
		
		We have proved that $E[1]\in\mathfrak C$. Solving the equation $\mathrm{ch}(E[1])=a\cdot\mathrm{ch}(\mathcal O_X(-1)[3])+b\cdot\mathrm{ch}(\mathcal Q(-1)[2])+c\cdot\mathrm{ch}(\mathcal U[1])+d\cdot\mathrm{ch}(\mathcal O_X)$ we find $a=-5e$, but $a$ should be a non-negative integer, hence $e\le 0$. For $e=0$ we find $b=3,c=6,d=1$, and Remark \ref{complexes} gives (\ref{3-term}).
	\end{proof}

\end{document}
